\section{Variation of the curve} 
\label{sectproofvariat}
%----------------------------------------------------------------------
In this appendix, we proove the theorems stated in sections (\ref{sectcompstruct})
and section (\ref{sectvarikappa}).

{\bf Lemma \ref{lemDOmega} }\proof{
\beq
\begin{array}{l}
D_\Omega\, \left( \sum_j \Res_{q\to a_j} {dE_{q}(p)\over \om(q)}\, f(q,\qbar) \right)_{x(p)}
\cr
=   \sum_j \Res_{q\to a_j} {dE_{q}(p)\over \om(q)} \,\, D_\Omega \left(f(q,\qbar)\right)_{x(q)} -  \sum_j \Res_{q\to a_j} {dE_{q}(p)\over (\om(q))^2}\,(\Omega(q)-\Omega(\qbar))\, f(q,\qbar)  \cr
 +2   \sum_j \Res_{q\to a_j} \sum_i \Res_{r\to a_i} { dE_{r}(p) \over \om(r)}\, \Omega(r)\, {dE_{q}(r) \over \om(q)}\, f(q,\qbar) \cr
 
=   \sum_j \Res_{q\to a_j} {dE_{q}(p)\over \om(q)} \,\, D_\Omega \left(f(q,\qbar)\right)_{x(q)} - 2  \sum_j \Res_{q\to a_j} {dE_{q}(p)\over (\om(q))^2}\,\Omega(q)\, f(q,\qbar)  \cr
 +2   \sum_j \Res_{q\to a_j} \sum_i \Res_{r\to a_i} { dE_{r}(p) \over \om(r)}\, \Omega(r)\, {dE_{q}(r) \over \om(q)}\, f(q,\qbar) \cr
=   \sum_j \Res_{q\to a_j} {dE_{q}(p)\over \om(q)} \,\, D_\Omega \left(f(q,\qbar)\right)_{x(q)} - 2  \sum_j \Res_{q\to a_j} \Res_{r\to q} {dE_{q}(r) dE_r(p)\over \om(q) \om(r)}\,\Omega(r)\, f(q,\qbar)  \cr
 +2   \sum_j \Res_{q\to a_j} \sum_i \Res_{r\to a_i} { dE_{r}(p) \over \om(r)}\, \Omega(r)\, {dE_{q}(r) \over \om(q)}\, f(q,\qbar) \cr
%
=   \sum_j \Res_{q\to a_j} {dE_{q}(p)\over \om(q)} \,\, D_\Omega \left(f(q,\qbar)\right)_{x(q)} \cr
 +2    \sum_i \Res_{r\to a_i}\sum_j \Res_{q\to a_j} { dE_{r}(p) \over \om(r)}\, \Omega(r)\, {dE_{q}(r) \over \om(q)}\, f(q,\qbar) \cr
\end{array}
\eeq
}

\vs

%%%%%%%%%%%%%%%%%%%%%%%%%%%%%%%%%%%%%%%%%%%%%%%%%%%%%%%%%%%%%%%%%%%%%%%%%%%%%
{\bf Theorem \ref{variat}} \proof{
This theorem straightforwardly comes from the diagrammatic rules described in the preceding paragraph except for
the variation of $F^{(1)}$.

Let us prove it for $F^{(1)}$ with $
\Omega(p) = \int_{\cal C} B(p,q) \Lambda(q)$.
One has
\beq
\begin{array}{rcl}
- \int_{\cal C} W_1^{(g)}(p) \L(p) &=&
- \Res_{q \to {\bf a}} {\int_{\cal C} dE_q(p) \L(p) \over \omega(q)} B(q,\qbar) \cr
&=& - \Res_{q \to {\bf a}} {\int_{\cal C} dE_q(p) \L(p) dz_i(q) dz_i(\qbar) \over \omega(q)} \left[ {1 \over (z(q)-z(\qbar))^2} + {1 \over 6} S_B(q) \right] \cr
\end{array}
\eeq
where $z_i$ is a local variable near the branch point $a_i$ and $S_B$ is the corresponding Bergmann projective connection.
Since the last term has a simple pole at the branch point $a_i$, one can write
\bea
- \Res_{q \to {\bf a}} {\int_{\cal C} dE_q(p) \L(p) dz_i(q) dz_i(\qbar) \over \omega(q)} {1 \over 6} S_B(q)
&=& - {1 \over 2} \sum_i {\Omega(a_i) \over dy(a_i)}  \Res_{q \to a_i} {B(q,\qbar) \over dx(q)} \cr
&=& - {1 \over 2} \delta_{\Omega} ln \tau_{Bx}. \cr
\eea

On the other hand, one can express the first term thanks to the local variables $z_i$ and compute
\bea
- \Res_{q \to {\bf a}} {\int_{\cal C} dE_q(p) \L(p) dz_i(q) dz_i(\qbar) \over \omega(q) (z(q)-z(\qbar))^2} 
&=& - {1 \over 24} \delta_\Omega (y'(a_i))\cr
\eea
provided that $z_i(q) - z_i(\qbar)= 2 z_i(q)$.
}

\vs

%%%%%%%%%%%%%%%%%%%%%%%%%%%%%%%%%%%%%%%%%%%%%%%%%%%%%%%%%%%%%%%%%%%%%%%%%%%
{\bf Theorem (\ref{homogeneity}) }\proof{
Using theorem \ref{thintPhi}, we have:
\beq
(2-2g)F^{(g)} = -\Res_\bfa \Phi W_1^{(g)}
\eeq
and using eq.\ref{ydxmoduli}, we can choose for any arbitrary $o'$:
\beq
\Phi(p)=
-\sum_\alpha \Res_{\alpha} V_\alpha dS_{p,o'}
+ \sum_\alpha t_\alpha \int_{o}^\alpha dS_{p,o'}
+ \sum_i \epsilon_i \oint_{\bcycle_i} dS_{p,o'}
\eeq
which implies:
\beq
\begin{array}{l}
 (2-2g)F^{(g)}
= -\Res_\bfa \Phi W_1^{(g)} \cr
= \sum_\alpha \Res_{p\to\bfa}  \Res_{q\to\alpha} V_\alpha(q) dS_{p,o'}(q) W_1^{(g)}(p)  - \sum_\alpha t_\alpha \Res_{p\to\bfa}  \int_{q=o}^\alpha dS_{p,o'}(q) W_1^{(g)}(p) \cr
 - \sum_i \epsilon_i \Res_{p\to\bfa} \oint_{q\in\bcycle_i} dS_{p,o'}(q) W_1^{(g)}(p) .\cr
\end{array}
\eeq
Since the poles $\alpha$ and the branch points $a_i$ do not coincide, one can exchange the order of integration. Then, one can
move the integration contours for $p$ in order to integrate only around the last pole $p \to q$:
\beq
\begin{array}{l}
 (2-2g)F^{(g)}
\cr
= \sum_\alpha t_\alpha  \int_{q=o}^\alpha \Res_{p\to q}  dS_{p,o'}(q) W_1^{(g)}(p) - \sum_\alpha  \Res_{q\to\alpha} \Res_{p\to q}  V_\alpha(q) dS_{p,o'}(q) W_1^{(g)}(p)   \cr
 + \sum_i \epsilon_i \oint_{q\in\bcycle_i} \Res_{p\to q}  dS_{p,o'}(q) W_1^{(g)}(p) \cr
=  \sum_\alpha  \Res_{q\to\alpha} V_\alpha(q) W_1^{(g)}(q)  - \sum_\alpha t_\alpha  \int_{q=o}^\alpha W_1^{(g)}(q)  - \sum_i \epsilon_i \oint_{q\in\bcycle_i} W_1^{(g)}(q) .\cr
\end{array}
\eeq
}


\vs

%%%%%%%%%%%%%%%%%%%%%%%%%%%%%%%%%%%%%%%%%%%%%%%%%%%%%%%%%%%
{\bf Theorem (\ref{thdWdkappa})} \proof{
We have
\bea
2i\pi {\partial \over \partial \kappa_{ij}}\,W^{(0)}_{2}(p_1,p_2)
&=& {1\over 2}(2i\pi)^2 (du_i(p_1) du_j(p_2)+du_i(p_2) du_j(p_1))\cr
&=& {1\over 2}\,\oint_{r\in\bcycle_j}\oint_{s\in\bcycle_i} (B(p_1,r)B(p_2,s) + B(p_2,r)B(p_1,s)). \cr
\eea
Thus the theorem holds for $k=2$ and $g=0$.
And by integration we get 
\bea
- 2i\pi {\partial \over \partial \kappa_{ij}} dE_q(p)
&=&  2 (i\pi)^2 (du_i(p) (u_j(q)-u_j(\qbar))+du_j(p) (u_i(q)-u_i(\qbar))) \cr
&=& - {1\over 2}\,\oint_{r\in \bcycle_i}\oint_{s\in \bcycle_j}  (B(p,r) dE_q(s)+B(p,s) dE_q(r)) .\cr
\eea

\beq
\begin{array}{rl}
& 2i\pi {\partial \over \partial \kappa_{ij}} W_{k+1}^{(g)}(p,p_K) \cr
%
=& 2i\pi {\partial \over \partial \kappa_{ij}} \Res_{q\to\bfa} {dE_q(p)\over (y(q)-y(\qbar))dx(q)}\, \Big(
W_{k+2}^{(g-1)}(q,\qbar,p_K) + \sum W_{j+1}^{(h)}(q,p_J) W_{k-j+1}^{(g-h)}(\qbar,p_{K/J}) \Big) \cr
%
=& {1 \over 2}\,\oint_{r\in \bcycle_i}\oint_{s\in \bcycle_j}  \Res_{q\to\bfa} {B(p,r) dE_q(s)+B(p,s) dE_q(r)\over (y(q)-y(\qbar))dx(q)}\, \Big(
W_{k+2}^{(g-1)}(q,\qbar,p_K) \cr
& \hspace{8cm} + \sum W_{j+1}^{(h)}(q,p_J) W_{k-j+1}^{(g-h)}(\qbar,p_{K/J}) \Big) \cr
&  + \Res_{q\to\bfa} {dE_q(p)\over (y(q)-y(\qbar))dx(q)}\, \Big( 2i\pi {\partial \over \partial \kappa_{ij}} W_{k+2}^{(g-1)}(q,\qbar,p_K) \Big) \cr
&+ \Res_{q\to\bfa} {dE_q(p)\over (y(q)-y(\qbar))dx(q)}\, \Big( \sum 2i\pi {\partial \over \partial \kappa_{ij}} W_{j+1}^{(h)}(q,p_J) W_{k-j+1}^{(g-h)}(\qbar,p_{K/J}) \Big) \cr
& + \Res_{q\to\bfa} {dE_q(p)\over (y(q)-y(\qbar))dx(q)}\, \Big( \sum  W_{j+1}^{(h)}(q,p_J)  2i\pi {\partial \over \partial \kappa_{ij}}W_{k-j+1}^{(g-h)}(\qbar,p_{K/J}) \Big) \cr
%
=&   {1 \over 2}\,\oint_{r\in \bcycle_i}\oint_{s\in \bcycle_j}  \Big( B(p,r) W^{(g)}_{k+1}(s,p_K)+B(p,s)  W^{(g)}_{k+1}(r,p_K) \Big) \cr
& + {1\over 2}\,\oint_{r\in \bcycle_i}\oint_{s\in \bcycle_j} \Res_{q\to\bfa} {dE_q(p)\over (y(q)-y(\qbar))dx(q)}\,
\Big( W_{k+4}^{(g-2)}(q,\qbar,p_K,r,s) \cr
& + 2 W_{j+3}^{(h)}(q,\qbar,p_J,r) W_{k-j+1}^{(g-1-h)}(p_{K/J},s) + 2 W_{j+2}^{(h)}(q,p_J,r) W_{k-j+2}^{(g-1-h)}(\qbar,p_{K/J},s)
\Big) \cr
%
& + {1\over 2}\,\oint_{r\in \bcycle_i}\oint_{s\in \bcycle_j} \Res_{q\to\bfa} {dE_q(p)\over (y(q)-y(\qbar))dx(q)}\,  \sum_J W_{k-j+1}^{(g-h)}(\qbar,p_{K/J}) \times \cr
&\hspace{2cm} \times  \Big( W_{j+3}^{(h-1)}(q,p_J,r,s) + 2 \sum_L W_{l+2}^{(m)}(q,p_L,r) W_{k-j-l+1}^{(h-m)}(p_{K/(J\cup L)},s) \Big) \cr
%
& + {1\over 2}\,\oint_{r\in \bcycle_i}\oint_{s\in \bcycle_j} \Res_{q\to\bfa} {dE_q(p)\over (y(q)-y(\qbar))dx(q)}\,  \sum_J W_{k-j+1}^{(g-h)}(q,p_{K/J}) \times \cr
& \hspace{2cm} \times \Big( W_{j+3}^{(h-1)}(\qbar,p_J,r,s) + 2 \sum_L W_{l+2}^{(m)}(\qbar,p_L,r) W_{k-j-l+1}^{(h-m)}(p_{K/(J\cup L)},s) \Big)   \cr
\end{array}
\eeq
We regroup together all terms with 2 $W$'s and all terms with 3 $W$'s:
\beq
\begin{array}{rl}
& 2i\pi {\partial \over \partial \kappa_{ij}} W_{k+1}^{(g)}(p,p_K) \cr
%
=&   {1 \over 2}\,\oint_{r\in \bcycle_i}\oint_{s\in \bcycle_j}  \Big( B(p,r) W^{(g)}_{k+1}(s,p_K)+B(p,s)  W^{(g)}_{k+1}(r,p_K) \Big) \cr
%
& + \oint_{r\in \bcycle_i}\oint_{s\in \bcycle_j} \Res_{q\to\bfa} {dE_q(p)\over (y(q)-y(\qbar))dx(q)}\,  \sum_J  W_{j+3}^{(h-1)}(q,p_J,r,s) W_{k-j+1}^{(g-h)}(\qbar,p_{K/J}) \cr
%
& + \oint_{r\in \bcycle_i}\oint_{s\in \bcycle_j} \Res_{q\to\bfa} {dE_q(p)\over (y(q)-y(\qbar))dx(q)}\,   \Big(  W_{j+2}^{(h)}(q,p_J,r) W_{k-j+2}^{(g-1-h)}(\qbar,p_{K/J},s) \Big) \cr
%
& + {1\over 2}\,\oint_{r\in \bcycle_i}\oint_{s\in \bcycle_j} \Res_{q\to\bfa} {dE_q(p)\over (y(q)-y(\qbar))dx(q)}\, W_{k+4}^{(g-2)}(q,\qbar,p_K,r,s) \cr
%
& + \oint_{r\in \bcycle_i}\oint_{s\in \bcycle_j} \Res_{q\to\bfa} {dE_q(p)\over (y(q)-y(\qbar))dx(q)} \times \cr
& \hspace{2cm} \times \Big( \sum_J  \sum_L W_{l+2}^{(m)}(\qbar,p_L,r) W_{k-j-l+1}^{(h-m)}(p_{K/(J\cup L)},s)  W_{k-j+1}^{(g-h)}(q,p_{K/J}) \Big) \cr
%
& + \oint_{r\in \bcycle_i}\oint_{s\in \bcycle_j} \Res_{q\to\bfa} {dE_q(p)\over (y(q)-y(\qbar))dx(q)} \times \cr
& \hspace{2cm} \times \Big( \sum_J   \sum_L W_{l+2}^{(m)}(q,p_L,r) W_{k-j-l+1}^{(h-m)}(p_{K/(J\cup L)},s)  W_{k-j+1}^{(g-h)}(\qbar,p_{K/J}) \Big) \cr
%
& + {1\over 2}\,\oint_{r\in \bcycle_i}\oint_{s\in \bcycle_j} \Res_{q\to\bfa} {dE_q(p)\over (y(q)-y(\qbar))dx(q)}\,  \Big( 2 W_{j+3}^{(h)}(q,\qbar,p_J,r) W_{k-j+1}^{(g-1-h)}(p_{K/J},s)  \Big) \cr
\end{array}
\eeq
We recognize the recursion  relation  eq.\ref{defWkgrecursive} in lines 2 to 5, and in lines 6 to 8, and this gives the theorem.

The theorem for the free energies, is easy obtained using theorem  \ref{thintPhi}.
}
